% CPEN 442 -- Assignment 4 answer template (optional).
% Fill in your name(s) and answers below, compile, and submit the PDF as a4.pdf.
% You do not have to use LaTeX: any readable PDF is fine.
\documentclass[12pt]{article}

\usepackage[margin=1in]{geometry}
\usepackage[dvipsnames,table]{xcolor}
\usepackage[most]{tcolorbox}
\usepackage{hyperref}
\usepackage{listings}
\usepackage{graphicx}
\usepackage{amsmath}
\usepackage{tabularx}
\usepackage[T1]{fontenc}
\usepackage{needspace}

\definecolor{alicepink}{rgb}{0.98,0.51,1.}
\definecolor{aliceblue}{rgb}{0.7, 1, 1}
\definecolor{background}{rgb}{1, 0.98, 0.95}

\lstset{frame=tb, language=Python, showstringspaces=false, columns=flexible,
  basicstyle={\small\ttfamily}, breaklines=true, backgroundcolor=\color{background}}

\setlength{\parindent}{0in}
\setlength{\parskip}{\baselineskip}

\newcommand{\question}[2]{\textbf{\colorbox{aliceblue}{[#2 pts]} Question #1}}

% Box to write each answer in. It grows with its content.
\newtcolorbox{answer}{colback=white, colframe=black, boxrule=0.6pt,
  sharp corners, breakable, height from=5cm to 50cm}

\begin{document}

\begin{center}
    THE UNIVERSITY OF BRITISH COLUMBIA\\
    Electrical and Computer Engineering Department\\[\baselineskip]
    {\bf CPEN 442\hfill Introduction to Cybersecurity\hfill Winter Term 1, 2026}
    \vspace{1cm}

    {\sc \Large Assignment 4: Cryptography Applications -- Searchable Encryption}
\end{center}

% One or two students. Delete the second row if you worked alone.
\begin{tabularx}{\textwidth}{|X|l|}
    \hline
    \rowcolor{gray!15} \textbf{Name} & \textbf{Student number} \\
    \hline
    First Last & 12345678 \\
    \hline
    First Last & 12345678 \\
    \hline
\end{tabularx}

The code for every question is in \texttt{a4.py}; this PDF explains how each attack works.
Write the explanations so that a classmate could follow them: in your own words, without unexplained jargon, and matching what your code actually does.

\section{Version 1: Hashing the Keywords}

\question{1}{15} Explain how \texttt{attack\_unkeyed} works.

\begin{answer}
    \textbf{Idea.}
    In Version 1, Alice turns a keyword $w$ into its query token with $t_w = h(\texttt{"token"} \,\|\, w)$.
    This formula contains no secret, so the server can evaluate it too.
    A hash cannot be run backwards, but I do not need to: the server knows the keyword space (all 200 possible keywords), so it can run the formula forwards on every candidate and compare the results with the tokens it received.

    \textbf{Steps.}
    \begin{enumerate}
        \item Build a lookup table: for each of the 200 keywords $w$ in the keyword space, compute $h(\texttt{"token"} \,\|\, w)$ and store \texttt{token $\to$ keyword}.
        \item For each query the server saw, take its token and look it up in the table.
        The keyword stored next to it is the answer for that query.
        \item Return the answers in the same order as the queries.
    \end{enumerate}

    \textbf{Example.}
    Say the keyword space is \{invoice, meeting, energy, price, contract\}.
    The table has five rows, e.g.\ \texttt{0c30e2f0...} $\to$ price.
    If Alice searches for price, the server sees the token \texttt{0c30e2f0...}, finds it in the table, and answers ``price''.
    The server never needs to look at the document IDs that came back.

    \textbf{Why it is always right.}
    The hash function is deterministic, so a keyword always gives the same token.
    It is also (in practice) collision-free, so no two keywords share a token.
    Hence the token of a query matches exactly one keyword, the one Alice searched for, and the attack recovers every query (QRR $=100\%$).
    It runs in time proportional to the size of the keyword space plus the number of queries.

    \textbf{Why the hashing does not help Alice.}
    One-wayness only protects secrets that are hard to guess.
    Here the set of possible inputs is tiny and public, so trying all of them is cheap, like a dictionary attack on a password hash.
    Version 2 fixes this by putting a secret key inside the token, which the server cannot compute.
\end{answer}

\Needspace{9cm}
\section{Version 2: The Secure Scheme}

\question{2}{75} Explain how \texttt{attack\_secure} works.

\begin{answer}
    \textbf{What the server can use.}
    The query tokens are now secret, but each query still shows which of Alice's $n$ documents matched.
    Two things can be read off this: (a) \emph{how many} documents a query returned, and (b) \emph{how many documents two queries have in common}.
    The same two numbers can be measured for every keyword on the auxiliary documents: how many contain the keyword, and how many contain a given pair of keywords.
    Since the auxiliary documents are similar to Alice's, a keyword's numbers there should be close to its numbers in Alice's documents.
    I also use that the 100 queries are for 100 \emph{different} keywords, so no keyword can be the answer for two queries.

    \textbf{Measuring a mismatch.}
    Suppose a query returned $k$ of Alice's $n$ documents, and a keyword occurs in a fraction $p$ of the auxiliary documents.
    I score how badly they fit with the binomial deviance
    \[ k\ln\frac{k}{np} + (n-k)\ln\frac{n-k}{n(1-p)}, \]
    which is $0$ when $k = np$ exactly and grows the further $k$ is from what $p$ predicts.
    (It is the negative log-probability of seeing $k$ matches, up to a constant, so it treats a small difference in a rare keyword as more important than in a common one.)
    The same score is used for the overlap of two queries, compared with the aux fraction of documents containing both keywords.

    \textbf{Step 1: guess from counts only.}
    For every (query, keyword) pair, the cost is the mismatch between the query's result count and the keyword's frequency in the aux data.
    I then pick the cheapest \emph{one-to-one} pairing of queries to keywords (the Hungarian algorithm, \texttt{linear\_sum\_assignment}), so no keyword is used twice.
    On its own this is weak, because many keywords have similar frequencies (about 12\% of queries correct on the Enron data).

    \textbf{Step 2: improve the guesses with overlaps.}
    Now use the current guesses of the \emph{other} queries.
    For query $i$ and a candidate keyword $w$, look at every other query $j$ with its current guess $v$.
    If $i$ really were $w$, the number of documents shared by $i$ and $j$ should be close to $n \cdot C[w][v]$, the aux co-occurrence of $w$ and $v$.
    The cost of ``$i$ is $w$'' is the count mismatch from Step 1 plus the sum of these overlap mismatches over all $j$.
    I skip a $j$ whose current guess is $w$ itself, since a keyword's overlap with itself is just its frequency, not a co-occurrence.
    Then I solve the one-to-one pairing again with the new costs.
    Better guesses give better overlap evidence, which gives better guesses, so I repeat this until the pairing stops changing (at most 30 rounds).

    \textbf{Example.}
    Take the hospital toy: if one query returned 1 of 6 documents it could be asthma, diabetes or fracture by count alone.
    But it shares no document with the query already identified as Covid, and in the aux data only fracture patients almost never have Covid, so the overlap cost is lowest for fracture.

    \textbf{Results and limits.}
    On the Enron data, averaged over 10 random splits, the attack recovers about 95\% of the queries in about 0.3 seconds per trial.
    On a separate set of 30 splits I did not tune on, it gets about 91\%, and it does worse when the aux data or Alice's collection is small (about 64\% and 72\% with a quarter of each).
    It can fail in an occasional split: if the first guesses are poor, the repeated updates can settle on a wrong pairing.
    Keywords that look the same to the aux data (similar frequency and similar co-occurrence) cannot be told apart at all.
\end{answer}

\Needspace{9cm}
\section{File Injection}

\question{3}{60} Explain how \texttt{choose\_injected\_documents} and \texttt{attack\_injection} work.

\begin{answer}
    \textbf{Idea.}
    Each document I inject tells me, for every query, a yes/no fact: ``does the queried keyword appear in my document?''
    With three documents, every keyword gets a 3-bit code (which of my documents contain it), so there are 8 groups of keywords.
    If I can read a query's code from its results, its keyword is one of only about 25 (instead of 200), and the attack of Question 2 has a much easier job.
    The difficulty is that Alice shuffles her documents, so I first have to find my own three documents among the thousands in the results.

    \textbf{\texttt{choose\_injected\_documents}.}
    I sort the keywords by how often they occur in the auxiliary documents and deal them out over the 8 codes like cards: the first keyword gets code 0, the second code 1, \dots, the ninth code 0 again.
    Document $j$ contains every keyword whose code has bit $j$ set, so each injected document holds half of the keyword space.
    Dealing by frequency puts keywords of similar frequency, the ones Question 2 mixes up most, into different groups.
    It also means my documents match queries equally often whether the keyword is common or rare, which a normal document does not.

    \textbf{\texttt{attack\_injection}, step 1: find my documents.}
    I do not know which of Alice's $n+3$ columns are mine, or which is which.
    Queries are a random half of the keywords, so a document with $s$ keywords should match about $s/2$ queries (give or take a few); columns far from that are ruled out.
    I then collect candidate triples of columns in two independent ways:
    \begin{enumerate}
        \item \emph{Agreement.} Run the plain attack of Question 2 (ignoring the injection), then for each of my documents look for the column that matches a query exactly when the guessed keyword is in that document.
        \item \emph{Looks like my documents.} Prefer columns whose matches do not depend on how common the query's keyword is (measured by how many results the query had), and triples that split the queries into the 8 code groups of roughly the expected sizes (about 12 queries each).
    \end{enumerate}
    Neither is reliable alone (the first fails if the plain attack happens to be poor, the second if Alice's documents are long and look like mine), so I decide between all candidate triples, in every order, by running the matching of Step 3 briefly for each and keeping the one with the lowest cost.
    A wrong triple gives wrong codes, so its cost is much higher.

    \textbf{Step 2: read each query's code.}
    A query's code is its pattern over my three columns: which of them it returned.
    It may only be assigned a keyword with that same code.

    \textbf{Step 3: match.}
    I remove my three columns, so counts and overlaps describe Alice's real documents (as in Question 2), and run the same one-to-one matching based on counts and overlaps with the aux data, but with each query allowed only the keywords of its group.

    \textbf{Example.}
    In the hospital toy, inject doc A = \{Covid, flu\} and doc B = \{flu, asthma\}.
    A query that returns both A and B must be flu, one that returns only B is asthma, one that returns only A is Covid, and one that returns neither is diabetes or fracture, which the counts and overlaps then separate.

    \textbf{Results and limits.}
    On the Enron data the attack recovers about 99\% of the queries (10 splits; 98.8\% and 99.0\% on two sets of 30 splits), compared with about 95\% without injection, in about 2 seconds.
    When I made Alice's documents three times longer, or gave myself a quarter of the aux data or of Alice's documents, it still found my documents in every trial and recovered 92--94\% (the attack without injection falls to 66--74\% there).
    The remaining mistakes are keywords inside the same group that the aux data cannot tell apart.
    It assumes the queries are a uniformly random subset of the keywords, which the assignment states.
\end{answer}

% Include this section only if you used generative AI.
% It must be between half a page and one page long, and cannot be written by AI.
\section*{Generative AI Disclosure}

% Describe how and where you used generative AI, and how helpful it was.

\end{document}
